% Biên dịch bằng XeLaTeX để hỗ trợ tiếng Việt Unicode.
\usepackage{fontspec}
\setmainfont{TeX Gyre Termes}
\usepackage{polyglossia}
\setmainlanguage{vietnamese}
\setotherlanguage{english}
\usepackage{graphicx}
\usepackage{geometry}
\usepackage{tocloft,calc}
\usepackage{amsmath,amsfonts,latexsym}
\usepackage{fancyhdr}
\usepackage{tikz}
\usepackage{eso-pic}
\usepackage{indentfirst}
\usepackage{setspace}
\usepackage[unicode,bookmarks=true,colorlinks=true]{hyperref}
\usepackage[capitalize]{cleveref}
% cleveref chưa có bộ tên ngôn ngữ Vietnamese; dùng hook tiếng Anh của nó,
% còn nhãn tham chiếu tiếng Việt được đặt ngay bên dưới.
\makeatletter
\def\cref@language{english}
\makeatother
\crefname{figure}{Hình}{Hình}
\Crefname{figure}{Hình}{Hình}
\crefname{table}{Bảng}{Bảng}
\Crefname{table}{Bảng}{Bảng}
\crefname{chapter}{Chương}{Chương}
\Crefname{chapter}{Chương}{Chương}
\usepackage{booktabs}
\usepackage{algorithm}
\usepackage{algpseudocode}
\floatname{algorithm}{Thuật toán}
\algrenewcommand\algorithmicrequire{\textbf{Đầu vào:}}
\algrenewcommand\algorithmicensure{\textbf{Đầu ra:}}
\usepackage{tabularx}
\usepackage{longtable}
\usepackage{multirow}
\usepackage[official]{eurosym}   % biểu tượng euro
\usepackage{enumitem}
\usepackage[table]{xcolor}
\usepackage{float}
\usepackage{chngcntr}
\usepackage[version=4]{mhchem}
\usepackage{fancybox}
\counterwithout{figure}{chapter}

\hypersetup{
  pdftitle={Khảo sát giải mã suy đoán không sử dụng bộ nhớ trọng số bổ sung trên GPU giới hạn tài nguyên},
  pdfauthor={Nguyễn Quốc Hiếu},
  linkcolor=blue,
  citecolor=blue,
  urlcolor=blue
}

\onehalfspacing

\geometry{
    top=3.5cm, 
    bottom=3.0cm, 
    left=3.5cm, 
    right=2.0cm
  % left=64pt,
  % right=64pt,
  % top=64pt,
  % bottom=64pt,
}

\usetikzlibrary{calc}

\newcommand{\drawcoverframe}{
    \AddToShipoutPicture*{%
    \begin{tikzpicture}[remember picture,overlay]
        % Khung ngoài
        \draw[line width=2pt]
            ($(current page.north west)+(3.5cm,-3.5cm)$)
            rectangle
            ($(current page.south east)+(-2cm,3cm)$);
    
        % Khung trong
        \draw[line width=1pt]
            ($(current page.north west)+(3.65cm,-3.65cm)$)
            rectangle
            ($(current page.south east)+(-2.15cm,3.15cm)$);
    \end{tikzpicture}
    }
}

\setlength{\headheight}{16pt}
\pagestyle{fancy}
\fancyhf{}
\lhead{}
\rhead{}
\cfoot{\thepage}

% Hide the printed "Chương <n>" heading label, but keep chapter numbering internally
\makeatletter
\def\@makechapterhead#1{%
  \vspace*{50\p@}%
  {\parindent \z@ \raggedright \normalfont
    \interlinepenalty\@M
    \Huge\bfseries #1\par\nobreak
    \vskip 40\p@
  }}
\def\@makeschapterhead#1{%
  \vspace*{50\p@}%
  {\parindent \z@ \raggedright \normalfont
    \interlinepenalty\@M
    \Huge\bfseries #1\par\nobreak
    \vskip 40\p@
  }}
\makeatother
